\chapter{Conclusions}

The development of a hybrid propulsion system through retrofitting involves multiple variables, with the technical and economic analysis of the battery bank as the central focus. The characterization of the load profile is the starting point, allowing the inefficiencies of the original system to be quantified and a battery bank to be sized according to the ship's operational requirements.\\
In addition to the technical analysis of the battery bank, it is crucial to integrate regulatory frameworks and physical limits (weight and volume), which are determining factors in the architecture of a tugboat. At the same time, profitability analysis—considering the optimal level of hybridization according to the asset's useful life—is vital to ensure return on investment (ROI) and attract capital interested in maritime decarbonization.\\

Once the technical and economic feasibility has been validated, mathematical modeling and subsequent analysis of the electrical performance of the battery bank and power converters are carried out. This comprehensive approach not only ensures that the system meets the demands of the propulsion system, but also establishes a robust methodology with multiple criteria, seeking to ensure that hybridization processes are successful and aligned with global emissions policies. In this way, this work seeks to lay the foundations for the creation of a methodology that allows the modeling of battery banks for hybrid tugboats, considering multiple variables in the design process.

